\chapter{Родительское происхождение спектральной меры нелокального ядра}
\label{ch:v8-baryon-nonlocal-kernel-spectral-measure-parent-origin-gate}

Предыдущий гейт допустил положительный нелокальный класс, но оставил его
спектральную меру свободной. Проверим сильнейший минимальный родитель: одно
вещественное вспомогательное поле и один положительный кинетический полюс.

Для составного трёхчастичного источника $J$ рассмотрим точную квадратичную
форму при $z=k^2\geq0$:
\begin{equation}
 S_{m,g}[\phi,J]
 =\frac12(z+m^2)\phi^2-g\phi J,
 \qquad m^2>0,quad g^2>0.
 \label{eq:v8-baryon-spectral-parent-action}
\end{equation}
Её вспомогательный гессиан положителен, а стационарная точка равна
\begin{equation}
 \phi_*(z)=\frac{g}{z+m^2}J.
 \label{eq:v8-baryon-spectral-parent-stationary-field}
\end{equation}
Точное исключение по дополнению Шура даёт
\begin{equation}
 S_{m,g}[\phi_*,J]
 =-\frac12\frac{g^2}{z+m^2}J^2.
 \label{eq:v8-baryon-spectral-parent-schur-kernel}
\end{equation}
Следовательно, одно вспомогательное поле действительно реализует
однополюсное нелокальное ядро.

Чтобы его статическая проекция совпадала с $\lambda_3W_3$, достаточно
\begin{equation}
 g^2=\lambda_3m^2,
 \qquad
 f_m(z)=\frac{m^2}{z+m^2},
 \qquad f_m(0)=1.
 \label{eq:v8-baryon-spectral-parent-static-matching}
\end{equation}
Но это условие фиксирует лишь отношение $g^2/m^2$. Для любого $q>0$
допустимо преобразование
\begin{equation}
 m^2\longmapsto q m^2,
 \qquad
 g^2\longmapsto q g^2.
 \label{eq:v8-baryon-spectral-parent-scale-orbit}
\end{equation}
Оно сохраняет статический коэффициент и все внутренние симметрии:
\begin{equation}
 \frac{qg^2}{qm^2}=\lambda_3,
 \qquad
 f_{q,m}(0)=1,
 \qquad
 f_{q,m}(z)=\frac{qm^2}{z+qm^2}.
 \label{eq:v8-baryon-spectral-parent-scale-invariance}
\end{equation}
Вне нуля формы различаются. Точные свидетели:
\begin{equation}
 f_{q,m}'(0)=-\frac1{qm^2},
 \qquad
 f_{1,m}(m^2)=\frac12,
 \qquad
 f_{2,m}(m^2)=\frac23.
 \label{eq:v8-baryon-spectral-parent-scale-witnesses}
\end{equation}

Эта свобода совпадает с уже найденной орбитой геометрии базы. Для
\begin{equation}
 \mathcal D_r=rD_M\otimes I_{21}+\gamma^5\otimes D_F,
 \qquad
 [\mathcal D_r,I_4\otimes a]=\gamma^5\otimes[D_F,a],
 \label{eq:v8-baryon-spectral-parent-base-scale-recall}
\end{equation}
все $42$ внутренних коммутатора независимы от $r$, тогда как внешний
квадрат импульса масштабируется. Поэтому конечный родитель, задающий $W_3$,
не различает членов одной спектральной орбиты.

\begin{theorem}[Минимальная реализация и масштабная неединственность]
Одно положительное вспомогательное поле реализует однополюсное нелокальное
ядро. После фиксации его статической проекции остаётся по меньшей мере
одномерная орбита $(m^2,g^2)\mapsto(qm^2,qg^2)$, на которой статический
оператор неизменен, а импульсная форма различна.
\end{theorem}

\begin{nogocriterion}[Конечный родитель не выбирает спектральную меру]
Нельзя считать массу полюса выведенной из $W_3$, его симметрий или условия
$K(0)=\lambda_3W_3$. Уже одно-полюсный родитель сохраняет свободный масштаб,
а прежняя конечная геометрия не выбирает масштаб оператора базы. Для
закрытия нужны независимый якорь геометрии базы либо новый спектральный
закон, различающий точки орбиты.
\end{nogocriterion}

\begin{protocol}
\begin{enumerate}
 \item поле вычислений: \textbf{$\mathbb Q(z,q,m^2,\lambda_3)$};
 \item чисел с плавающей точкой: \textbf{$0$};
 \item вспомогательных полей: \textbf{$1$};
 \item положительный гессиан поля: \textbf{да};
 \item дополнение Шура: \textbf{точное};
 \item статическое условие: \textbf{$g^2/m^2=\lambda_3$};
 \item размерность остаточной орбиты: \textbf{не меньше $1$};
 \item спектральная масса выбрана: \textbf{нет};
 \item родительское происхождение меры: \textbf{NO-GO};
 \item следующий гейт: \textbf{аудит кандидатов якоря спектрального масштаба}.
\end{enumerate}
\end{protocol}

Машинный сертификат сохранён в
\path{s2t_v8_baryon_nonlocal_kernel_spectral_measure_parent_origin_gate_results.json}.